% Documento escrito final (paper) - Proyecto final MAT 4021 Matemáticas Discretas
% Compilar:  tectonic paper.tex      (o pdflatex/latexmk con una distribución TeX)
\documentclass[conference,a4paper]{IEEEtran}

\usepackage{iftex}
\ifTUTeX
  % XeLaTeX / LuaLaTeX (Tectonic usa XeTeX)
  \usepackage{fontspec}
  \setmainfont{texgyretermes}[Extension=.otf, UprightFont=*-regular,
    BoldFont=*-bold, ItalicFont=*-italic, BoldItalicFont=*-bolditalic]
\else
  % pdfLaTeX
  \usepackage[utf8]{inputenc}
  \usepackage[T1]{fontenc}
  \usepackage{newtxtext}
\fi
\usepackage[spanish,es-tabla,es-noshorthands]{babel}
\usepackage{amsmath,amssymb,amsthm}
\usepackage{booktabs}
\usepackage{graphicx}
\usepackage{tikz}
\usetikzlibrary{positioning}
\usepackage{xcolor}
\usepackage[hidelinks]{hyperref}
\usepackage{url}

\newtheorem{definicion}{Definición}
\newtheorem{proposicion}{Proposición}

\newcommand{\repo}{\url{https://github.com/lau-240/Proyecto-final-Matem-ticas-Discretas}}
\newcommand{\F}{\mathcal{F}}

\begin{document}

\title{Disponibilidad y uso de espacios del campus:\\
un modelo con álgebra booleana y grafos ponderados}

\author{\IEEEauthorblockN{Laura Vanessa Acuña Ordóñez, Isabella Barrios Pico, Ángel Gabriel Carrillo Henao,\\
Montserrat Mariottis Martínez, Javier Enrique Torregrosa Lascano}
\IEEEauthorblockA{MAT 4021 Matemáticas Discretas -- Universidad del Norte\\
Barranquilla, Colombia\\
Repositorio: \repo}}

\maketitle

\begin{abstract}
Encontrar un salón libre para estudiar o reunirse es una tarea frecuente en el
campus, y suele resolverse caminando de bloque en bloque. Este trabajo propone un
modelo discreto para responder dos preguntas: \emph{qué espacios están libres en un
intervalo de tiempo} y \emph{cuál de ellos queda más cerca}. La ocupación semanal de
cada espacio se representa como un vector de 96 bits, uno por franja horaria. Las
consultas y la detección de conflictos se reducen a operaciones del álgebra
booleana $(\mathcal{P}(\F),\cup,\cap,\overline{\phantom{x}})$. Se distingue entre
\emph{libre} y \emph{sin información} mediante un segundo vector de bits, lo que da un
modelo de tres estados. El campus se modela como un grafo ponderado de bloques, en el
que el algoritmo de Dijkstra ordena los espacios libres por minutos de caminata y las
medidas de centralidad identifican los bloques más accesibles y los de paso. El
modelo se implementó en Python con una interfaz web en Streamlit y se validó con un
caso de resultado conocido, 104 pruebas automáticas y una simulación. En la
simulación, con 21\,\% de cobertura, el modelo ingenuo de dos estados se equivoca en
33\,\% de las respuestas ``libre''. El modelo de tres estados no responde ``libre''
sin información que lo respalde.
\end{abstract}

\begin{IEEEkeywords}
álgebra booleana, vectores de bits, grafos ponderados, Dijkstra, centralidad,
asignación de espacios.
\end{IEEEkeywords}

\section{Introducción}
En la universidad los horarios de clase ocupan los salones solo en ciertas franjas.
El resto del tiempo los espacios podrían usarse para estudiar, hacer tutorías o
reunirse, pero no existe un lugar donde consultar esa disponibilidad. Un estudiante
que busca un salón libre suele recorrer varios bloques sin éxito.

El problema tiene una estructura discreta natural. El tiempo se divide en franjas de
una hora, un horario es un subconjunto de franjas y el campus es un conjunto de
bloques unidos por caminos. Este trabajo aprovecha esa estructura con dos
herramientas del curso:
\begin{enumerate}
  \item \textbf{Álgebra booleana y conjuntos}: cada horario es un subconjunto del
  conjunto de franjas, codificado como un entero cuyos bits indican pertenencia.
  \item \textbf{Teoría de grafos}: los bloques forman un grafo no dirigido con pesos en
  minutos de caminata, sobre el que se calculan caminos mínimos y centralidades.
\end{enumerate}

La contribución principal es un modelo de \emph{tres estados} (ocupado, libre,
desconocido). Evita el error de suponer que un espacio sin información está libre.
El sistema se alimenta de los horarios que cargan los propios estudiantes, de modo que
la cobertura de la información crece con el uso.

\section{Marco teórico}

\subsection{Conjuntos y álgebra booleana}
Sea $\F$ un conjunto finito. El conjunto potencia $\mathcal{P}(\F)$ con la unión, la
intersección y el complemento forma un álgebra booleana con elemento cero
$\varnothing$ y elemento unidad $\F$~\cite{rosen}. Si $|\F|=n$, la función
\[
  \chi:\mathcal{P}(\F)\to\{0,1\}^n,\qquad \chi(S)_f = \begin{cases}1 & f\in S\\ 0 & f\notin S\end{cases}
\]
es un isomorfismo de álgebras booleanas entre $\mathcal{P}(\F)$ y $\{0,1\}^n$ con las
operaciones bit a bit OR, AND y NOT. Así, cada operación de conjuntos se calcula con
una sola instrucción sobre enteros: $S\cup T\mapsto s\,|\,t$, $S\cap T\mapsto s\,\&\,t$,
$\overline{S}\mapsto \lnot s\,\&\,U$ (donde $U=2^n-1$ representa a $\F$), y
$S\subseteq T \iff s\,\&\,t = s$.

\subsection{Grafos ponderados y caminos mínimos}
Un grafo ponderado es una terna $G=(V,E,w)$ con $w:E\to\mathbb{R}_{>0}$. La distancia
$d(u,v)$ es el peso mínimo de un camino entre $u$ y $v$. El algoritmo de
Dijkstra~\cite{dijkstra} calcula $d(s,\cdot)$ desde un vértice $s$ en
$O((|V|+|E|)\log|V|)$ con una cola de prioridad, siempre que los pesos sean no
negativos.

\subsection{Centralidad}
La \emph{cercanía} de un vértice $v$ es
$C(v)=\dfrac{|V|-1}{\sum_{u\neq v} d(u,v)}$, y es alta para los vértices desde los que
se llega rápido a todos los demás. La \emph{intermediación}
\[
  B(v)=\sum_{s\neq v\neq t}\frac{\sigma_{st}(v)}{\sigma_{st}},
\]
donde $\sigma_{st}$ es el número de caminos mínimos de $s$ a $t$ y $\sigma_{st}(v)$
los que pasan por $v$, mide cuánto actúa $v$ como lugar de paso~\cite{freeman}. En el
cálculo se usa normalizada.

\section{Modelo}

\subsection{Franjas horarias}
La semana se divide en franjas de una hora, de lunes a sábado, entre las 6:00 y las
22:00: $6\times16=96$ franjas. La franja del día $d\in\{0,\dots,5\}$ (0 = lunes) que
empieza a la hora $h$ es
\[
  f(d,h) = 16\,d + (h-6), \qquad \F=\{0,1,\dots,95\}.
\]
Un intervalo del día $d$ entre $h_1$ y $h_2$ (sin incluir $h_2$) es el conjunto
$W=\{f(d,h): h_1\le h<h_2\}$. Su máscara es un bloque de $h_2-h_1$ bits consecutivos,
$\chi(W) = (2^{h_2-h_1}-1)\cdot 2^{f(d,h_1)}$.

\subsection{Modelo de tres estados}
Para cada espacio $e$ se guardan dos subconjuntos de $\F$:
\begin{itemize}
  \item $O_e$: franjas en las que $e$ está \emph{ocupado} (clase, tutoría o uso
  registrado);
  \item $I_e$: franjas cuyo estado es \emph{conocido}.
\end{itemize}

\begin{definicion}
El estado de $e$ en la franja $f$ es
\[
\text{estado}(e,f)=\begin{cases}
\text{ocupado} & f\in O_e,\\
\text{libre} & f\in I_e\setminus O_e,\\
\text{desconocido} & f\notin I_e.
\end{cases}
\]
\end{definicion}

Toda ocupación registrada es también información, así que el sistema mantiene el
invariante $O_e\subseteq I_e$. Por tanto las tres clases forman una partición de $\F$.

\begin{definicion}[Disponibilidad en una ventana]
Un espacio $e$ está \emph{disponible} en la ventana $W\subseteq\F$ si
\[
  O_e\cap W=\varnothing \quad\text{y}\quad W\subseteq I_e,
\]
es decir, en bits: $o_e\,\&\,w = 0$ y $i_e\,\&\,w = w$.
\end{definicion}

Si el usuario lo pide de forma explícita, se relaja la segunda condición y se aceptan
espacios \emph{sin conflicto conocido}, aunque falte información. En la interfaz, una
ventana con parte de sus franjas conocidas y ninguna ocupada se muestra como
\emph{parcial}.

\subsection{Registro de usos y conflictos}
Registrar un uso de $e$ en $W$ se acepta si y solo si $O_e\cap W=\varnothing$. En ese
caso se actualiza $O_e\leftarrow O_e\cup W$ e $I_e\leftarrow I_e\cup W$, lo que
preserva el invariante. Si otro estudiante reporta el mismo uso (mismo espacio,
franjas, motivo y materia, aunque la escriba distinto), no se cuenta como conflicto
sino como una confirmación adicional.

\begin{proposicion}
El costo de decidir un conflicto o la disponibilidad de un espacio es $O(1)$
operaciones sobre enteros de 96 bits, independientemente de cuántas franjas tenga
$W$.
\end{proposicion}

Las clases y tutorías se repiten cada semana hasta el fin del semestre. Los usos de
una sola vez (estudio, reunión, práctica) vencen al terminar su próxima ocurrencia.
Al vencer, el sistema hace $O_e\leftarrow O_e\setminus W$ e
$I_e\leftarrow I_e\setminus W$: que alguien haya usado un salón una vez no prueba que
quede libre después.

\subsection{Grafo del campus}
Los vértices son los bloques del campus (A--M y el bloque de laboratorios LAB,
$|V|=14$). Las aristas unen bloques vecinos y su peso son los minutos de caminata
($|E|=17$, Fig.~\ref{fig:grafo}). Cada espacio pertenece a un bloque, así que la
distancia de un bloque origen $s$ a un espacio $e$ es $d(s,\text{bloque}(e))$.

\begin{figure}[t]
\centering
\begin{tikzpicture}[
  v/.style={circle,draw,fill=blue!10,minimum size=6.5mm,inner sep=0pt,font=\footnotesize},
  lab/.style={font=\scriptsize,fill=white,inner sep=1pt}, scale=1]
  \foreach \n/\x in {A/0,B/1.2,C/2.4,D/3.6,E/4.8,F/6}
    \node[v] (\n) at (\x,0) {\n};
  \foreach \n/\x in {L/0,K/1.2,J/2.4,I/3.6,H/4.8,G/6}
    \node[v] (\n) at (\x,-2) {\n};
  \node[v] (M) at (0,-3.4) {M};
  \node[v,font=\tiny] (LAB) at (1.4,-3.4) {LAB};
  \draw (A)--node[lab]{2}(B); \draw (B)--node[lab]{2}(C);
  \draw (C)--node[lab]{3}(D); \draw (D)--node[lab]{2}(E);
  \draw (E)--node[lab]{2}(F); \draw (F)--node[lab]{3}(G);
  \draw (G)--node[lab]{2}(H); \draw (H)--node[lab]{2}(I);
  \draw (I)--node[lab]{3}(J); \draw (J)--node[lab]{2}(K);
  \draw (K)--node[lab]{2}(L); \draw (L)--node[lab]{6}(M);
  \draw (L)--node[lab]{3}(LAB);
  \draw (B) to[bend left=35] node[lab]{4} (E);
  \draw (C) to[bend left=35] node[lab]{5} (F);
  \draw (F)--node[lab]{5}(I);
  \draw (H) to[bend left=30] node[lab]{4} (K);
\end{tikzpicture}
\caption{Grafo de bloques del campus. Los pesos son minutos de caminata de ejemplo
(\texttt{ARISTAS\_CAMPUS} en \texttt{src/datos.py}).}
\label{fig:grafo}
\end{figure}

\subsection{Espacio libre más cercano}
Dada una ventana $W$, un bloque origen $s$ y opcionalmente un tipo de espacio (salón,
laboratorio, cancha, piscina, sala de música), el sistema:
\begin{enumerate}
  \item calcula $d(s,\cdot)$ con Dijkstra;
  \item filtra los espacios disponibles en $W$ con las operaciones de bits;
  \item ordena el resultado por $(d(s,\text{bloque}(e)),\ \text{id}(e))$.
\end{enumerate}
El costo total es $O((|V|+|E|)\log|V| + n\log n)$ para $n$ espacios.

Si el estudiante cargó su horario, el origen $s$ se infiere. Es el bloque de la clase
en curso; si no hay ninguna, el de la última clase que terminó antes; si tampoco la
hay, el de la próxima clase. Sin horario, el estudiante indica el bloque a mano.

\section{Implementación}
El sistema está escrito en Python. Usa NetworkX~\cite{networkx} para el grafo y
Streamlit para la interfaz web. Los módulos principales son:
\begin{itemize}
  \item \texttt{franjas.py}: franjas como bits (máscaras, rangos, conteo).
  \item \texttt{campus.py}: espacios, vectores $O_e$ e $I_e$, registros, conflictos y
  grafo.
  \item \texttt{consultas.py}: disponibilidad, ranking por cercanía, tutorías,
  ubicación inferida y cobertura.
  \item \texttt{codigos.py}: traduce códigos de salón (\texttt{26J}, \texttt{37G2},
  \texttt{SDU9}, \texttt{L1-5}) a su bloque.
  \item \texttt{almacen.py}: guardado en JSON por semestre, con candado de escritura.
  \item \texttt{lector\_imagen.py}: lectura opcional de la foto del horario con un
  modelo de IA. El resultado siempre lo revisa el estudiante antes de guardarlo.
\end{itemize}

La interfaz tiene seis pestañas: \emph{Buscar libre}, \emph{¿Qué hay?},
\emph{Tutorías}, \emph{Registrar}, \emph{Mi horario} y \emph{Semestre}. Cada
estudiante se identifica con un código inventado, no con su nombre. Los horarios vencen
con el semestre y el administrador puede archivarlos o borrarlos.

\section{Validación y resultados}

\subsection{Caso de resultado conocido}
Se construyó un campus pequeño, con tres bloques y aristas $A\!-\!B$ (2 min) y
$B\!-\!C$ (3 min). Tiene cuatro salones con estados conocidos el lunes de 10:00 a
12:00: A1 ocupado por una tutoría, A2 sin información, y B1 y C1 libres y
confirmados. La estudiante tiene clase en A1 de 8:00 a 10:00, así que a las 10:00 su
ubicación inferida es el bloque A. El resultado esperado y el obtenido coinciden:
\begin{itemize}
  \item solo con datos confirmados: \textbf{B1 a 2 min};
  \item aceptando desconocidos: \textbf{A2 a 0 min}, marcado como desconocido;
  \item registrar B1 de 14 a 16 y luego de 15 a 17 da conflicto, porque
  $W_1\cap W_2=\{15\text{:}00\}\neq\varnothing$.
\end{itemize}

\subsection{Pruebas automáticas}
El repositorio incluye 104 pruebas con \texttt{pytest}. Cubren franjas, conflictos,
vencimiento de usos, códigos de salón, lectura de horarios, almacenamiento
concurrente y la interfaz. Se ejecutan en cada \emph{push} con GitHub Actions. Las 104
pasan.

\subsection{Simulación: tres estados frente a dos}
Se generó una ocupación ``real'' aleatoria para un campus de 42 espacios. Cada bloque
de dos horas está ocupado con probabilidad $0{,}45$. Luego se revela cada par
(espacio, franja) con probabilidad igual a la \emph{cobertura}. Se compara el modelo
de tres estados ($\text{libre}=I\cap\overline{O}$) con un modelo ingenuo de dos
estados ($\text{libre}=\overline{O}$, sin información $=$ libre). La precisión es la
fracción de respuestas ``libre'' que son correctas.

\begin{table}[h]
\centering
\caption{Precisión de la respuesta ``libre'' según la cobertura.}
\label{tab:sim}
\begin{tabular}{@{}ccc@{}}
\toprule
Cobertura & Modelo 3 estados & Modelo ingenuo \\
\midrule
0,21 & 1,00 & 0,67 \\
0,41 & 1,00 & 0,73 \\
0,61 & 1,00 & 0,80 \\
0,80 & 1,00 & 0,89 \\
0,95 & 1,00 & 0,98 \\
\bottomrule
\end{tabular}
\end{table}

Con poca información (21\,\%), una de cada tres respuestas ``libre'' del modelo
ingenuo es falsa. La precisión de 1,00 del modelo de tres estados se da \emph{por
construcción}, porque la simulación supone que los datos confirmados son correctos.
En la práctica el riesgo está en las confirmaciones erróneas, y por eso el sistema
cuenta cuántos estudiantes reportan lo mismo. El costo del modelo de tres estados es
la \emph{exhaustividad}: con poca cobertura encuentra menos espacios. Por eso la
interfaz permite incluir los desconocidos, marcados como tales.

\subsection{Análisis de la red}
Sobre el grafo de la Fig.~\ref{fig:grafo}, los bloques con mayor cercanía son G y H
($C=0{,}137$) y F ($0{,}134$). Son los mejores lugares para ubicar servicios comunes.
Los de mayor intermediación son F ($B=0{,}513$), G y H ($0{,}385$). Por ellos pasa la
mayor parte de los recorridos más cortos, y un cierre en F alargaría muchos caminos.

\section{Limitaciones}
\begin{itemize}
  \item Los pesos del grafo son valores de ejemplo y deben medirse en el campus.
  \item Un espacio solo se conoce si alguien lo registró o cargó un horario que lo
  usa. ``Sin información'' no significa ``libre''.
  \item El código estudiantil no es secreto. Quien conozca el de otra persona puede
  modificar su horario.
  \item La persistencia en JSON con candado sirve para un solo proceso. Con muchos
  usuarios simultáneos convendría una base de datos.
\end{itemize}

\section{Conclusiones}
El álgebra booleana sobre vectores de bits da una representación compacta y exacta de
los horarios. Con ella, las preguntas de disponibilidad y conflicto se reducen a
operaciones de tiempo constante. Separar ``ocupado'' de ``conocido'' convierte un
modelo ingenuo, que puede equivocarse en un tercio de los casos, en uno que no afirma
nada sin respaldo. El grafo ponderado y Dijkstra convierten la lista de espacios
libres en una recomendación útil (el más cercano), y las centralidades describen la
estructura del campus. Como trabajo futuro quedan medir los tiempos reales de
caminata, cargar horarios reales y comparar, con una encuesta, el tiempo de búsqueda
manual frente al del modelo.

\section*{Código fuente}
Todo el código, las pruebas y las instrucciones para ejecutarlo están en:
\begin{center}\footnotesize\repo\end{center}

\begin{thebibliography}{9}
\bibitem{rosen} K.~H. Rosen, \emph{Matemática discreta y sus aplicaciones}, 7.\textordfeminine{} ed.
Madrid: McGraw-Hill, 2012.
\bibitem{dijkstra} E.~W. Dijkstra, ``A note on two problems in connexion with
graphs,'' \emph{Numerische Mathematik}, vol.~1, pp.~269--271, 1959.
\bibitem{freeman} L.~C. Freeman, ``A set of measures of centrality based on
betweenness,'' \emph{Sociometry}, vol.~40, no.~1, pp.~35--41, 1977.
\bibitem{networkx} A.~A. Hagberg, D.~A. Schult y P.~J. Swart, ``Exploring network
structure, dynamics, and function using NetworkX,'' en \emph{Proc. 7th Python in
Science Conference (SciPy)}, 2008, pp.~11--15.
\end{thebibliography}

\end{document}
